\subsection{Substitution in the functional calculus}

With the notations introduced above, the statements of Cor. 1 of I, p.~66 and Cor. 2 of I, p.~67 become respectively

\[
\chi (g (\boldsymbol {a})) = g \left(\chi \left(a _ {1}\right), \dots , \chi \left(a _ {n}\right)\right), \quad \operatorname{Sp} (g (\boldsymbol {a})) = g \left(\operatorname{Sp} ^ {n} (\boldsymbol {a})\right)
\]

for \(f \in \mathcal{O}(\mathrm{Sp}^n(\boldsymbol{a});\mathrm{A})\) , \(\chi \in \mathrm{X}(\mathrm{A})\) and \(g \in \mathcal{O}(\mathrm{Sp}^n(\boldsymbol{a}))\) .

We shall now prove a more general substitution property.

THEOREM 4. --- Let A be a commutative unital complex Banach algebra, let \(n \geqslant 1\) be an integer and \(\boldsymbol{a} = (a_1, \ldots, a_n) \in \mathrm{A}^n\) . Let \(\boldsymbol{f} = (f_1, \ldots, f_p)\) where \(f_1, \ldots, f_p\) are elements of \(\mathcal{O}(\mathrm{Sp}^n(\boldsymbol{a}))\) . The image of \(\mathrm{Sp}^n(\boldsymbol{a})\) by the map \(\boldsymbol{z} \mapsto \boldsymbol{f}(\boldsymbol{z}) = (f_1(\boldsymbol{z}), \ldots, f_p(\boldsymbol{z}))\) is equal to \(\mathrm{Sp}^p(\boldsymbol{f}(\boldsymbol{a}))\) .

For every \(g \in \mathcal{O}(\mathrm{Sp}^{p}(\boldsymbol{f}(\boldsymbol{a}));\mathrm{A})\) , the composed germ \(g \circ f\) is an element of \(\mathcal{O}(\mathrm{Sp}^{n}(\boldsymbol{a});\mathrm{A})\) and we have \(g(\boldsymbol{f}(\boldsymbol{a})) = (g \circ \boldsymbol{f})(\boldsymbol{a})\) .

The first assertion concerning the image of \(\mathrm{Sp}^{n}(\boldsymbol{a})\) follows from Cor. 2 of I, p.~67. To prove the second, we shall use the following lemma.

Lemma 12. --- Let K be the polynomially convex hull of \(\mathrm{Sp}^p (\pmb {f}(\pmb {a}))\) . We have \(g(\pmb {f}(\pmb {a})) = (g\circ \pmb {f})(\pmb {a})\) for every germ \(g\in \mathcal{O}(\mathrm{K};\mathrm{A})\) .

Let \(\Psi\) be the map from \(\mathcal{O}(\mathrm{K};\mathrm{A})\) to A such that \(\Psi(g) = (g\circ f)(\boldsymbol{a})\) . It is a continuous unital morphism, such that \(\Psi(z_j) = f_j(\boldsymbol{a})\) , where \(z_{j}\) is the germ of the \(j\) -th coordinate function on \(\mathbf{C}^p\). When \(g\) is the germ of a polynomial function, we therefore have \(\Psi(g) = g(\boldsymbol{f}(\boldsymbol{a}))\). By th. 2 of I, p.~68, this formula remains valid for all \(g\in \mathcal{O}(\mathrm{K};\mathrm{A})\) .

Now let us prove the theorem. Let V be an open neighborhood of \(\mathrm{Sp}^{p}(\boldsymbol{f}(\boldsymbol{a}))\) and \(\widetilde{g}\in\mathcal{O}(V;A)\) be a holomorphic function whose germ in a neighborhood of \(\mathrm{Sp}^{p}(\boldsymbol{f}(\boldsymbol{a}))\) is equal to g. Let \(b\in C^{p+q}\) be an envelope of \((\boldsymbol{f}(\boldsymbol{a}),\mathrm{V})\) (Lemma 11 of I, p.~70) and \(\pi\colon V\times C^{q}\to V\) the canonical projection.

Let \(\widetilde{f}_1, \ldots, \widetilde{f}_p\) be holomorphic functions whose germs are \(f_1, \ldots, f_p\) and let U be an open neighborhood of \(\mathrm{Sp}^n(\boldsymbol{a})\) such that \((\widetilde{f}_1, \ldots, \widetilde{f}_p)(\mathrm{U}) \subset \mathrm{V}\) . Let \(\pi'\) be the canonical projection from \(\mathrm{U} \times \mathbf{C}^q\) onto U. Denote by \(z_{n+1}, \ldots, z_{n+q}\) the \(q\) last coordinate functions on \(\mathbf{C}^{n+q}\) . Denote by \(h = \widetilde{g} \circ (\widetilde{f}_1, \ldots, \widetilde{f}_p)\) and

\[
\boldsymbol {c} = (a _ {1}, \dots , a _ {n}, b _ {p + 1}, \dots , b _ {p + q}) \in \mathrm{A} ^ {n + q}.
\]

The map \(\widetilde{g} \circ \pi\) is holomorphic in the open neighborhood \(V \times C^q\) of the polynomially convex hull L of \(Sp^{p+q}(b)\). By Lemma 12, applied to \(c\) and to the germs in the neighborhood of L of the functions

\[
(\widetilde {f} _ {1} \circ \pi^ {\prime}, \ldots , \widetilde {f} _ {p} \circ \pi^ {\prime}, z _ {n + 1}, \ldots , z _ {n + q}),
\]

and to the germ of \(\widetilde{g}\circ\pi\) , one has

\[
(g \circ \pi) \left((f _ {1} \circ \pi^ {\prime}) (\boldsymbol {c}), \dots , (f _ {p} \circ \pi^ {\prime}) (\boldsymbol {c}), z _ {n + 1} (\boldsymbol {c}), \dots , z _ {n + q} (\boldsymbol {c})\right) = (h \circ \pi^ {\prime}) (\boldsymbol {c}).
\]

Since \(\pi'(\boldsymbol{c}) = \boldsymbol{a}\) , one has \((h \circ \pi')(\boldsymbol{c}) = h(\boldsymbol{a})\) and \((f_i \circ \pi')(\boldsymbol{c}) = f_i(\boldsymbol{a})\) for \(1 \leqslant i \leqslant p\) (property (CF3) of the holomorphic functional calculus). Moreover, since \(z_{n+j}(\boldsymbol{c}) = b_{p+j}\) for \(1 \leqslant j \leqslant q\). By property (CF2), we have

\[
(g \circ \pi) (f _ {1} (\boldsymbol {a}), \dots , f _ {p} (\boldsymbol {a}), b _ {p + 1}, \dots , b _ {p + q}) = h (\boldsymbol {a}),
\]

from which we deduce \(g(f_{1}(\boldsymbol{a}), \ldots, f_{p}(\boldsymbol{a})) = h(\boldsymbol{a})\) by applying property (CF3) again.

